\section{Introduction}
\label{sec:intro}

World models that predict the future in the feature space of a frozen vision foundation model are cheap to train and plan with \citep{zhou2025dinowm,baldassarre2025dinoworld,assran2025vjepa2}, but a deterministic predictor returns the conditional mean, which lies between the modes of a multimodal future \citep{porcher2026flowwm}. Stochastic feature-space world models fix this with diffusion or flow matching at tens of network evaluations (NFE) per step \citep{porcher2026flowwm,boduljak2025vfmf,walker2025generalist}, a cost that matters when a planner imagines thousands of futures per decision. Several recent works therefore build one-step stochastic world models \citep{lu2026driftworld,agibot2026geact,lee2026beckmann,park2026priming,hsu2026vsf}.

A simpler one-step objective predates these methods: draw $m$ candidate futures from a noise-conditioned generator, find the one closest to the observed outcome, and pull only that one towards it. This objective is conditional implicit maximum likelihood estimation \citep[IMLE;][]{li2018imle,li2019diverse}; it is also the min-of-$K$ objective of explorative modeling \citep[XM;][]{gladstone2026xm}, whose $K$ is the budget we call $m$, and a sampled relative of winner-takes-all learning \citep{guzman2012mcl,rupprecht2017mhp}. Because every data point is explained by at least one sample, mode collapse is discouraged by the objective itself, and inference is a single forward pass. IMLE world models beat Gaussian ones on proprioceptive state \citep{aghabozorgi2026wimle}, and \citet{gladstone2026xm} conjecture that feature-space world models have few modes, so that a small candidate budget should suffice. In the opposite direction, \citet{makabe2026esp} note that a hard minimum over finitely many candidates is not a strictly proper scoring rule and show that an energy-score generator covers a two-branch toy distribution better than min-of-$K$.

This paper asks a quantitative version of both claims: \emph{how many candidates does a one-step \nom{} world model need, when is no budget enough, and is a strictly proper score the better objective?} We answer with a conditional IMLE world model on frozen DINOv2/v3 patch tokens, compared at matched parameter counts against deterministic regression, flow matching with 50 and 2 steps in the FlowWM recipe, two one-step flow generators (improved MeanFlow and Shortcut models) and an energy-score generator, on benchmarks whose futures can be enumerated, so that coverage is measured rather than inferred. One quantity organizes the results: the \nom{} loss is, in expectation, the distortion of quantizing the outcome with a \emph{random codebook} of $m$ draws from the generator. This is a sampled version of the known equivalence between winner-takes-all learning and conditional quantization \citep{rupprecht2017mhp,letzelter2024wta,perera2024amcl}, and in the continuous limit it coincides with the Gaussian rate--distortion test channel \citep{shannon1959coding,cover2006elements}. The mathematics is classical; our contribution is to show that this single correspondence predicts, in advance of the runs, where the budget matters and where it cannot help.

\paragraph{Contributions.}
(1) \textbf{A controlled test of the small-budget conjecture.} Under mode-like uncertainty, \nom{} with $m{=}16$ covers enumerated futures better than 50-step flow matching (FM-50), one-step flow generators and an energy-score generator at matched parameter counts (coverage $0.93$ vs.\ $0.79$ and $0.36$ on non-repeating contexts; F1 error $7.2$ vs.\ $12.3$\,px on Bouncing Shapes) and at $30\times$ lower latency than FM-50; on a stochastic Push-T variant its precision advantage is significant and its recall advantage consistent but not significant (Section~\ref{sec:results1}).
(2) \textbf{Outcome memorization and its incentive.} In contexts that never repeat, a conditional generator can reach zero training loss by memorizing the realized outcome of each training context, keyed on its appearance. We measure this as the share of samples on the realized mode of training contexts, show that it is governed by the codebook's miss \emph{distortion} $I(m)$ rather than by the miss probability $(1-1/K)^m$ (fit RMSE $0.031$ vs.\ $0.068$ over $K \in \{2,4,8,16\}$), confirm seven of nine predictions recorded before the $K{=}2,16$ runs, and explain why the strictly proper energy score memorizes at every $m$ while \nom{} with $m{=}16$ does not (Section~\ref{sec:results2}).
(3) \textbf{Where no budget suffices.} Mode probabilities flatten with $m$, matching the optimum of \citet{bennett2026freedom} within $0.01$, and the spread on continuous high-dimensional outcomes shrinks to $e_m/\sqrt{d}$ of the true spread, the form of the Gaussian rate--distortion test channel; on BAIR the samples spread along only one to three context-specific directions, and $m{=}64$ does not close the gap to multi-step flow matching on distribution metrics (Section~\ref{sec:results3}).
(4) \textbf{Rollouts and planning.} \nom{} rollouts stay as close to the token manifold as deterministic regression, while flow rollouts drift by up to $4\times$ at $h{=}32$ regardless of the number of integration steps; per-step matching on self-generated contexts \citep[in the lineage of][]{venkatraman2015dad,zhang2025catk} cuts long-horizon recall error by a third on stochastic Push-T, and an ablation separates its manifold benefit from its coverage benefit; planning with \imlewm{} is $50\times$ faster than with FM-50 at equal success, but success does not improve even where the model's risk estimates are accurate (Section~\ref{sec:results4}).

Negative results are reported in Appendix~\ref{app:negative}. The benchmarks are synthetic or simulated except for BAIR, the models are small (8M and 21M parameters), and the flow baselines are far from the scale of \citet{porcher2026flowwm}; Section~\ref{sec:discussion} states these limits.
